\section{S-integers, residue symbols, and Gauss sums}\label{sec:arithmetic}

The analytic argument will sum over ideals while also using their generators
in local fields.  We first choose generators for which reciprocity has no
supplementary factor.  We then identify the two Hecke characters which account
for the phases of the cubic and sextic Gauss sums.

Fix a cyclotomic field \(F\) containing \(\mu_{12}\), and put
\([F:\Q]=2r\).  At a finite place \(v\), normalize the absolute value by
\(|\varpi_v|_v=(\Norm v)^{-1}\); at a complex place use the square of the
ordinary modulus.  These are the absolute values in the product formula.
Let
\[
 \psi=\prod_v\psi_v:\Adele_F/F\longrightarrow \C^\times
\]
be the standard trace additive character.  We also fix a finite-order Hecke
character \(\eta\).  Choose a finite set \(S\) containing the complex places,
all places above \(6\), the finite conductor of \(\eta\), and every place at
which \(\mathcal O_v\) is not self-dual for the pairing
\(\psi_v(xy)\).  The theta normalization in \Cref{prop:theta-input}
and the local Euler-product estimate in \Cref{prop:euler-product} require
two additional fixed finite sets of excluded places.  Include both sets in
\(S\).  Their existence is established at those points, independently of
the summation variables and of the arithmetic choices below.  We may
include any further fixed finite set of places.
Enlarging \(S\) by finitely many primes which generate the ideal class group,
we arrange that the ring of S-integers
\[
 R=\mathcal O_S=\{a\in F:v_{\mathfrak p}(a)\geq0
                       \text{ for every }\mathfrak p\notin S\}
\]
is principal.  This is the final choice of \(S\), and all arithmetic data
below are constructed for it before the coefficient family is defined.
All assertions in this section concern this fixed \(F\) and \(S\); their
auxiliary constants may depend on them.

Write \(S_f\) for the finite places in \(S\), \(s_f=|S_f|\), and
\[
 U_S=R^\times,\qquad F_S=\prod_{v\in S}F_v,\qquad
 |y|_S=\prod_{v\in S}|y_v|_v,\qquad
 F_S^1=\{y\in F_S^\times:|y|_S=1\}.
\]
An \emph{exterior ideal} is a nonzero integral ideal of \(R\), equivalently
an ideal supported at primes outside \(S\).  Its norm is
\(q_{\mathfrak a}=\Norm\mathfrak a=|R/\mathfrak a|\), extended
multiplicatively to fractional ideals.  Roman letters will denote generators
of the corresponding fraktur ideals.  For \(a\in F^\times\), put
\(q_a=|a|_S\).  Finally, on the additive group \(F_S\) set
\[
 \e(x)=\prod_{v\in S}\psi_v(x_v),
 \qquad
 \widehat f(y)=\int_{F_S}f(x)\e(-xy)\,dx,
\]
where \(dx\) is the product of the self-dual local additive measures.

\begin{lemma}[The S-integer lattice]\label{lem:s-lattice}
The diagonal copy of \(R\) is a discrete cocompact subgroup of \(F_S\), is
its own annihilator for the pairing \(\e(xy)\), and has covolume \(1\).
For every Bruhat--Schwartz function \(f\) on \(F_S\),
\[
 \sum_{m\in R}f(m)=\sum_{\nu\in R}\widehat f(\nu).
\]
If \(\mathfrak a=aR\) is an exterior fractional ideal, then
\[
 q_a=q_{\mathfrak a},\qquad
 (aR)^\perp=a^{-1}R,\qquad
 \sum_{m\in aR}f(m)
   =q_{\mathfrak a}^{-1}\sum_{\nu\in a^{-1}R}\widehat f(\nu).
\]
\end{lemma}

\begin{proof}
Put \(K^S=\prod_{v\notin S}\mathcal O_v\), embedded in the adeles with zero
S-components.  At every place outside \(S\),
the defining conductor condition says that \(\mathcal O_v\) is its own
annihilator.  Its self-dual measure is consequently \(1\), and
\(\widehat{\one_{K^S}}=\one_{K^S}\).  Adelic additive duality identifies the
annihilator of \(F\) in \(\Adele_F\) with \(F\), and the additive quotient
\(\Adele_F/F\) has volume \(1\); see
\cite[Chapter~II, \S29]{Weil64}.  Strong approximation identifies
\[
 F_S/R\simeq \Adele_F/(F+K^S).
\]
The quotient is compact because \(\Adele_F/F\) is compact.  To check
discreteness, choose a compact neighborhood \(C\) of zero in \(F_S\).
The intersection \(F\cap(C\times K^S)\) is finite because diagonal \(F\)
is discrete and closed in the adeles.  Shrinking \(C\) excludes the finitely
many nonzero S-components in this intersection, so zero is isolated in
\(R\).  Taking annihilators in the last quotient gives
\[
 (F+K^S)^\perp=F\cap(F_S\times K^S)=R,
\]
with \(R\) diagonally embedded, so its annihilator in \(F_S\) is indeed \(R\).
The covolume of a lattice and that of its annihilator are reciprocal under
self-dual measure; hence this covolume is \(1\).

More directly, apply adelic Poisson summation
\cite[Chapter~III, \S41, equation~(39)]{Weil64} to
\(f\otimes\one_{K^S}\).  Only the elements of \(R\) occur on either side,
and the exterior factor is unchanged by Fourier transform.  This gives the
displayed Poisson identity with no scalar factor.  Multiplication by \(a\)
has additive modulus \(|a|_S\), so scaling this identity gives the formula
for \(aR\).  The product formula gives
\[
 |a|_S=\prod_{v\notin S}|a|_v^{-1}
      =\Norm(aR)=q_{\mathfrak a},
\]
which finishes the proof.
\end{proof}

We use Hilbert symbols to remove the supplementary factor in sextic
reciprocity.  Choose the local reciprocity map \(\operatorname{Art}_v\) in
which a uniformizer acts as arithmetic Frobenius on an unramified extension.
For \(n\in\{2,3,6\}\), set
\[
 (a,b)_{n,v}
   =\frac{\operatorname{Art}_v(b)(a^{1/n})}{a^{1/n}}.
\]
This orders the two arguments oppositely to the convention in
\cite[Chapter~V, \S3]{Neukirch99}.  Thus, at a place of residue
characteristic prime to \(n\), for a unit \(u\) and a uniformizer \(\varpi\),
\[
 (u,\varpi)_{n,v}\equiv u^{(\Norm v-1)/n}\pmod v.
\]
The value is read in \(\mu_n\subset F\).  On
\[
 \mathcal V=F_S^\times/F_S^{\times6}
\]
define \(H_S(x,y)=\prod_{v\in S}(x_v,y_v)_{6,v}\), and let \(E\) be the image
of \(U_S\) in \(\mathcal V\).  For \(T>0\), let \(\ell_T\in F_S^\times\)
have finite components \(1\) and component \(T^{1/(2r)}\) at every complex
place.  Then \(|\ell_T|_S=T\).

\begin{lemma}[Generators with fixed power classes]\label{lem:balanced-generators}
The pairing \(H_S\) on \(\mathcal V\) is nondegenerate and alternating.
The map \(U_S/U_S^6\to\mathcal V\) is injective, its image \(E\) satisfies
\(E^\perp=E\), and there is an isotropic subgroup \(B\) such that
\(\mathcal V=E\oplus B\).

Every exterior fractional ideal \(\mathfrak a\) has a generator \(a\) whose
class \([a]\) in \(\mathcal V\) belongs to \(B\).  Any two such generators
differ by an element of \(U_S^6\).  They can be chosen so that
\[
 a/\ell_{q_{\mathfrak a}}\in\mathcal K_{\mathrm{bal}}
\]
for one fixed compact subset \(\mathcal K_{\mathrm{bal}}\subset F_S^1\), independent of
\(\mathfrak a\).
\end{lemma}

\begin{proof}
For \(\ell=2,3\), put \(\mathcal V_\ell=F_S^\times/F_S^{\times\ell}\).
Local Kummer theory and local reciprocity give a nondegenerate Hilbert
pairing on this group \cite[Chapter~V, Proposition~(3.2)]{Neukirch99}.
At a finite place \(v\),
\[
 \dim_{\Fq_\ell}F_v^\times/F_v^{\times\ell}
  =2+\one_{v\mid\ell}[F_v:\Q_\ell].
\]
Indeed the valuation contributes one dimension.  When \(v\nmid\ell\), the
residue units contribute one dimension and the principal units are
\(\ell\)-divisible.  When \(v\mid\ell\), the principal units have a
\(\Z_\ell\)-part of rank \([F_v:\Q_\ell]\) and an \(\ell\)-power
root-of-unity part contributing one more dimension.  The latter is present
because \(F\) contains \(\mu_\ell\).  This is also the local power-class
formula \cite[Chapter~II, Corollary~(5.8)]{Neukirch99}.
Complex power-class groups are trivial.
All places above \(\ell\) lie in \(S\), so their local degrees sum to \(2r\)
and
\begin{equation}
 \dim_{\Fq_\ell}\mathcal V_\ell=2s_f+2r.\label{eq:local-power-dimension}
\end{equation}
The S-unit theorem \cite[Chapter~I, Corollary~(11.7)]{Neukirch99} gives
\(U_S\simeq\mu(F)\times\Z^{r+s_f-1}\).  Since \(\mu(F)\) is cyclic and
contains \(\mu_\ell\),
\begin{equation}
 \dim_{\Fq_\ell}U_S/U_S^\ell=r+s_f.\label{eq:unit-power-dimension}
\end{equation}

We next check that the latter group injects into \(\mathcal V_\ell\).
Suppose that \(u\in U_S\) is an \(\ell\)-th power in every \(F_v\) with
\(v\in S\).  The Kummer extension \(F(u^{1/\ell})/F\) is cyclic and split
at \(S\).  At every exterior prime \(u\) is a unit and \(\ell\) is
invertible, so this extension is unramified there.  It is consequently an
everywhere unramified abelian extension in which every S-prime splits.
The Hilbert class-field correspondence, together with the quotient
description of the S-class group, identifies such extensions with quotients
of the class group of \(R\), which is trivial
\cite[Chapter~I, Proposition~(11.6); Chapter~VI, Proposition~(6.9) and
Theorem~(7.3)]{Neukirch99}.  Hence
\(u=t^\ell\) for \(t\in F\).  Its exterior valuations show \(t\in U_S\),
proving the injection.

For two S-units, every exterior Hilbert symbol is \(1\), since both arguments
are units and the residue characteristic is prime to \(6\).  Hilbert
reciprocity \cite[Chapter~VI, Theorem~(8.1)]{Neukirch99} therefore makes
their S-product \(1\).  The Hilbert pairing is
alternating here: \(-1=\zeta_{12}^6\), and the diagonal identity
\((x,x)_{n,v}=(x,-1)_{n,v}\) makes its diagonal trivial for \(n=2,3,6\).
By \Cref{eq:local-power-dimension} and \Cref{eq:unit-power-dimension}, the
injected unit image in \(\mathcal V_\ell\) is an isotropic subspace of half
the dimension.  Nondegeneracy implies that it is self-orthogonal.
Symplectic linear algebra supplies an isotropic complementary subspace for
each \(\ell\).  The Chinese remainder isomorphism
\(\mathcal V\simeq\mathcal V_2\oplus\mathcal V_3\), under which the two
pairings are the third and second powers of \(H_S\), combines these
complements into \(B\).  It also gives the asserted injection and
self-orthogonality for exponent \(6\).

Let \(\mathfrak a=a_0R\).  Decompose \([a_0]=e b\) in \(E\oplus B\), and
choose an S-unit representing \(e^{-1}\).  Its product with \(a_0\) is a
generator \(a\) with \([a]\in B\).  If \(a'\) is another, then \(a'/a\) is
an S-unit whose class lies in \(E\cap B\).  It is therefore locally a sixth
power at \(S\), and the injection already proved says \(a'/a\in U_S^6\).

Finally \(a/\ell_{q_{\mathfrak a}}\) belongs to \(F_S^1\), by
\Cref{lem:s-lattice}.  To see the relevant compactness explicitly, the
norm-one idele class group is compact
\cite[Chapter~VI, Theorem~(1.6)]{Neukirch99}.  The subgroup
\(F_S^1\times\prod_{v\notin S}\mathcal O_v^\times\) is open in the
norm-one ideles and meets diagonal \(F^\times\) in \(U_S\).  Its quotient
by \(U_S\) is therefore an open, hence closed, subgroup of that compact
idele class group.  Its projection onto \(F_S^1/U_S\) is surjective, so
\(F_S^1/U_S\) is compact.  Since \(U_S^6\) has finite index in \(U_S\),
\(F_S^1/U_S^6\) is compact as well.  Choosing a compact set which maps onto
this quotient and multiplying \(a\) by a suitable element of \(U_S^6\)
proves the last assertion.
\end{proof}

Call a generator with class in \(B\) \emph{allowed}, and call it
\emph{balanced} when it satisfies the compactness condition in
\Cref{lem:balanced-generators}.  We use balanced representatives for
estimates involving S-components.  In multiplicative identities we may
instead use the product of the allowed generators of the prime factors:
the product remains allowed and differs from the selected representative by
an element of \(U_S^6\).

For an exterior prime \(\mathfrak p\), the reduction of \(\mu_6\) into
\((R/\mathfrak p)^\times\) is injective.  Define the sextic residue character
by
\[
 \chi_{\mathfrak p}(x)\equiv x^{(q_{\mathfrak p}-1)/6}\pmod{\mathfrak p}
 \quad (x\not\equiv0\pmod{\mathfrak p}),\qquad
 \chi_{\mathfrak p}(0)=0.
\]
For every integer \(j\), the notation \(\chi_{\mathfrak p}(x)^j\) means the
usual character power on units and means \(0\) on nonunits, also when
\(j\equiv0\pmod6\).  If
\(\mathfrak a=\prod_{\mathfrak p}\mathfrak p^{k_{\mathfrak p}}\), put
\[
 \chi_{\mathfrak a}(x)=
 \prod_{\mathfrak p\mid\mathfrak a}
       \chi_{\mathfrak p}(x)^{k_{\mathfrak p}},
 \qquad \chi_R(x)=1.
\]
The same zero convention applies to powers of \(\chi_{\mathfrak a}\).

For distinct exterior primes with allowed generators \(p,p'\), global
reciprocity gives the exact equality
\begin{equation}
 \chi_{\mathfrak p}(p')=\chi_{\mathfrak p'}(p).
\label{eq:sextic-reciprocity}
\end{equation}
Indeed the only possibly nontrivial exterior symbols in the product formula
for \((p,p')_6\) are
\((p,p')_{6,\mathfrak p}=\chi_{\mathfrak p}(p')^{-1}\) and
\((p,p')_{6,\mathfrak p'}=\chi_{\mathfrak p'}(p)\).  The S-product is \(1\)
because \(B\) is isotropic.  If \(m\in R\setminus\{0\}\) and
\(\mathfrak b=(m)R\), write \(m=\epsilon b\)
with \(b\) allowed and \(\epsilon\in U_S\).  Multiplying
\Cref{eq:sextic-reciprocity} over prime factors gives, for an allowed
generator \(a\) of \(\mathfrak a\),
\begin{equation}
 \chi_{\mathfrak a}(m)
  =\chi_{\mathfrak a}(\epsilon)\chi_{\mathfrak b}(a).
\label{eq:generator-reciprocity}
\end{equation}
This includes the case \((\mathfrak a,\mathfrak b)\ne R\): both sides are
then zero.  For an S-unit \(\epsilon\), the same product formula gives
\begin{equation}
 \chi_{\mathfrak a}(\epsilon)=H_S(\epsilon,a)^{-1}.
\label{eq:unit-symbol}
\end{equation}
Thus its dependence on \(\mathfrak a\) is through the finite class \([a]\).

\begin{lemma}[Kummer and class characters]\label{lem:kummer-characters}
For \(u\in R\setminus\{0\}\), the ideal function
\[
 \chi_\bullet(u)(\mathfrak a)=\chi_{\mathfrak a}(u)
\]
is induced away from \(S\) and the prime divisors of \((u)R\) by a
finite-order Hecke character.  Its values at those prime divisors are set
equal to zero.  If \(\mathfrak f_u\) is the conductor of the inducing
primitive character, then
\begin{equation}
 \Norm\mathfrak f_u\le C_{F,S}\Norm\rad((u)R)
                       \le C_{F,S}q_u.\label{eq:kummer-conductor}
\end{equation}
The same bound, with a changed constant, holds after twisting by the fixed
\(\eta\).  The inducing characters associated to \(u\) and \(v\) agree if
and only if \(u/v\in F^{\times6}\).

Every character of the finite group \(B\), evaluated on allowed generators,
extends to a finite-order Hecke character unramified outside \(S\).

The following parametrization will index the nonzero Poisson frequencies;
its character distinctness and conductor bound will be used in zero detection.
Choose a complete set of representatives \(\mathcal U\) modulo \(U_S^6\) among the nonzero
S-integers whose exterior valuations belong to \(\{0,1,\ldots,5\}\), with
\(1\in\mathcal U\).  For allowed generators \(r\), the map
\begin{equation}
 (u,\mathfrak r)\longmapsto [u r^6]
 \quad\text{from }\mathcal U\times\{\text{exterior ideals}\}
 \quad\text{to }(R\setminus\{0\})/U_S^6\label{eq:frequency-decomposition}
\end{equation}
is a bijection.  Among \(u\in\mathcal U\), only \(u=1\) induces the
principal character, and distinct \(u\)'s induce distinct characters, also
after the fixed twist by \(\eta\).
\end{lemma}

\begin{proof}
Because \(\mu_6\subset F\), Kummer theory
\cite[Chapter~IV, Theorem~(3.3) and Corollary~(3.6)]{Neukirch99}
identifies \(u\in F^\times\) with
the character
\[
 \kappa_u(\sigma)=\frac{\sigma(u^{1/6})}{u^{1/6}}
 \quad\text{of }\Gal(\overline F/F),
\]
whose kernel as a function of \(u\) is exactly \(F^{\times6}\).
Global class-field theory identifies this finite abelian character with a
finite-order Hecke character
\cite[Chapter~VI, Theorem~(5.5) and Proposition~(5.6)]{Neukirch99}.
At an exterior prime
\(\mathfrak p\nmid(u)R\), arithmetic Frobenius satisfies
\[
 \kappa_u(\operatorname{Frob}_{\mathfrak p})
 \equiv (u^{1/6})^{q_{\mathfrak p}-1}
 =u^{(q_{\mathfrak p}-1)/6}
 =\chi_{\mathfrak p}(u)\pmod{\mathfrak p}.
\]
Reduction is injective on \(\mu_6\), so these are equal as roots of unity.
This proves the asserted identification on good primes.

At an exterior place, the principal units are sixth powers by Hensel's
lemma.  Hence the associated local Hilbert character is trivial on principal
units, giving conductor exponent at most one.  The tame Hilbert formula
\cite[Chapter~V, Proposition~(3.4)]{Neukirch99} shows that the pairing of
two units is \(1\), so if \(u\) is a local unit the character is unramified.
At a place in \(S\), it depends only on the class of \(u\) in
the finite group \(F_v^\times/F_v^{\times6}\); the finitely many resulting
local conductors have a fixed maximum.  This proves
\Cref{eq:kummer-conductor}.  A fixed twist supported in \(S\) only changes
that maximum.  Equality of two inducing Hecke characters is equality of
their Kummer characters, since unramified Frobenius elements determine a
finite abelian character.  The Kummer kernel gives the stated criterion.
Setting extra prime values to zero merely omits Euler factors and does not
alter the inducing character.

For a character \(\theta_B\) of \(B\), extend it to \(\mathcal V=E\oplus B\)
by \(\theta_S(eb)=\theta_B(b)\).  This is a continuous finite-order character
of \(F_S^\times\) trivial on \(U_S\).  If \(\mathfrak a=aR\), define
\(\theta(\mathfrak a)=\theta_S(a)\).  It is independent of the generator
and multiplicative.  At every exterior place take the unramified local
character sending a uniformizer to \(\theta(\mathfrak p)\), and at \(S\)
take \(\theta_S^{-1}\).  Their product is trivial on diagonal \(F^\times\)
by construction and is the required Hecke character.

To prove \Cref{eq:frequency-decomposition}, divide each exterior valuation
of a nonzero \(h\) by \(6\).  This determines an ideal \(\mathfrak r\) and
the exterior valuations of an S-integer \(u_0=h/r^6\) in
\(\{0,\ldots,5\}\).  Replacing \(u_0\) by its representative in
\(\mathcal U\) changes \(u_0r^6\) only by an element of \(U_S^6\).
Conversely, the exterior valuations of \([ur^6]\) determine
\(\mathfrak r\) and those of \(u\); equality of two such classes then says
that the two representatives \(u\) differ by \(U_S^6\).  Changing an
allowed \(r\) by \(U_S^6\) does not affect its image in the quotient.
This proves bijectivity.

If \(u\in\mathcal U\) induces the principal character, the first part gives
\(u=t^6\) in \(F\).  Its exterior valuations are divisible by \(6\) and
between \(0\) and \(5\), hence are all zero; the same is true for the
exterior valuations of \(t\).  Thus \(u\in U_S^6\), giving \(u=1\).
If \(u,v\in\mathcal U\) induce the same character, the differences of their
exterior valuations are divisible by \(6\) and lie in \([-5,5]\), so vanish.
Their quotient is then a sixth power of an S-unit, and \(u=v\).
The fixed twist by \(\eta\) cancels from an equality of twisted characters.
\end{proof}

We now determine the phases of the Gauss sums.  If \(\mathfrak c\) is a
squarefree exterior ideal with allowed generator \(c\), define
\begin{equation}
 \gamma_j(\mathfrak c)
  =q_{\mathfrak c}^{-1/2}
       \sum_{x\bmod\mathfrak c}\chi_{\mathfrak c}(x)^j\e(x/c),
 \qquad \gamma_j(R)=1.\label{eq:normalized-gauss}
\end{equation}
The additive character on \(R/\mathfrak c\) here is primitive by
\Cref{lem:s-lattice}.  A change \(c\mapsto\epsilon^6c\), with
\(\epsilon\in U_S\), changes the additive argument by a sixth power and is
absorbed by a change of residue variable; hence the sum is independent of
the allowed generator.  At an exterior prime, the adelic product formula
gives
\[
 \e(x/p)=\psi_{\mathfrak p}(-x/p).
\]
The reduction of \(\mu_{12}\) is injective there, so
\(q_{\mathfrak p}\equiv1\pmod{12}\) and
\(\chi_{\mathfrak p}(-1)=1\).  The minus sign can therefore be suppressed
in all these prime Gauss sums.

\begin{lemma}[The Gauss correction characters]\label{lem:gauss-characters}
There are a unitary Hecke character \(\Lambda\), unramified outside \(S\),
and a finite-order Hecke character \(G\), unramified outside \(S\), such that
for every exterior prime \(\mathfrak p\),
\begin{equation}
 \gamma_2(\mathfrak p)^3=-\Lambda(\mathfrak p),\qquad
 G(\mathfrak p)=\overline{\chi_{\mathfrak p}(4)}\,
                     \gamma_3(\mathfrak p),\qquad
 \gamma_1(\mathfrak p)\gamma_2(\mathfrak p)
                      =-\Lambda(\mathfrak p)G(\mathfrak p).
\label{eq:gauss-corrections}
\end{equation}
At each complex place the infinity type of \(\Lambda\) is purely angular
of exponent \(1\) or \(-1\), in the ordinary angular coordinate.
The values \(\gamma_3(\mathfrak p)\) are signs and depend multiplicatively
on the class of an allowed \(p\) in \(B\).

For coprime squarefree exterior ideals,
\begin{equation}
 \gamma_j(\mathfrak c_1\mathfrak c_2)
 =\chi_{\mathfrak c_1}(c_2)^j
  \chi_{\mathfrak c_2}(c_1)^j
  \gamma_j(\mathfrak c_1)\gamma_j(\mathfrak c_2).
\label{eq:gauss-crt}
\end{equation}
\end{lemma}

\begin{proof}
First work over a finite field \(k\) of size \(Q\), with a nontrivial
additive character \(\psi_k\), a character \(\chi\) of exact order \(6\), and
\(\rho=\chi^2\).  Write
\[
 G_j=\sum_{x\in k}\chi(x)^j\psi_k(x),\qquad
 J(\alpha,\beta)=\sum_{x\in k}\alpha(x)\beta(1-x).
\]
For nontrivial characters the usual Gauss inversion and Gauss--Jacobi
identities give
\begin{equation}
 G_2G_4=Q,\qquad G_2^2=J(\rho,\rho)G_4,\qquad
 (Q^{-1/2}G_2)^3=Q^{-1/2}J(\rho,\rho).\label{eq:cubic-jacobi}
\end{equation}
The first equality uses \(\rho(-1)=1\).  Changing \(\psi_k(x)\) to
\(\psi_k(ax)\) multiplies \(G_j\) by \(\overline{\chi(a)}^{\,j}\), so the
cube in \Cref{eq:cubic-jacobi} is independent of the additive character.

Let \(E_0=\Q(\omega)\), where \(\omega^2+\omega+1=0\).  Every ideal of
\(\Z[\omega]\) prime to \(3\) has a unique generator congruent to \(1\)
modulo \(3\); denote it by \(\pi_+(\mathfrak a)\).  This follows from class
number one and the bijection of the six units with
\((\Z[\omega]/3)^\times\), and is the convention used in
\cite[\S2.1]{BaierYoung10}.  For a prime \(\mathfrak q\nmid6\) with residue
size \(Q\), take
\(\rho(x)\equiv x^{(Q-1)/3}\pmod{\mathfrak q}\).  Reduction of the Jacobi
sum gives
\[
 J(\rho,\rho)\equiv
 \sum_{x\in\Fq_Q}x^{(Q-1)/3}(1-x)^{(Q-1)/3}=0
 \pmod{\mathfrak q},
\]
since every monomial has degree strictly between \(0\) and \(Q-1\).
The Gauss identities also give \(|J(\rho,\rho)|^2=Q\).
To determine the unit, put \(t=1-\omega\); the ideals \((t^2)\) and \((3)\)
are equal.  Writing \(\rho(x)=\omega^{k(x)}\), with \(k(x)\in\{0,1,2\}\),
and expanding modulo \(t^2\) yields
\[
 J(\rho,\rho)\equiv Q-2
   -t\sum_{x\ne0,1}\bigl(k(x)+k(1-x)\bigr)
 \equiv Q-2-2(Q-1)t\equiv-1\pmod3.
\]
Here each value of \(\rho\) occurs \((Q-1)/3\) times on
\(\Fq_Q^\times\), and \(Q\equiv1\pmod3\).  Divisibility by
\(\mathfrak q\), the norm equality, and this congruence show
\begin{equation}
 J(\rho,\rho)=-\pi_+(\mathfrak q).
\label{eq:eisenstein-jacobi}
\end{equation}
This is also the prime case of the classical cubic Gauss identity
\(g(n)^3=\mu(n)\Norm(n)n\) in
\cite[\S2.2, equation~(11) in arXiv:0804.2233v4]{BaierYoung10}.

Define on ideals of \(E_0\) prime to \(3\)
\[
 \Lambda_{E_0}(\mathfrak a)
     =\frac{\pi_+(\mathfrak a)}{|\pi_+(\mathfrak a)|}.
\]
The generators congruent to \(1\pmod3\) multiply, so this is a
multiplicative ideal character.  On a principal ideal generated by
\(a\equiv1\pmod3\) its value is \(a/|a|\); thus it is the unitary Hecke
character of finite conductor dividing \((3)\) and angular infinity type
of exponent \(1\), in ideal-character convention.  Set
\(\Lambda=\Lambda_{E_0}\circ\Norm_{F/E_0}\).
If \(\mathfrak p\) lies above \(\mathfrak q\) with residue degree \(f\),
the power definition of the residue character shows that the cubic
character at \(\mathfrak p\) is the lift of that at \(\mathfrak q\) by
the residue-field norm.  The Davenport--Hasse Jacobi lifting formula says
\[
 J_{\Fq_{Q^f}}(\rho\circ\Norm,\rho\circ\Norm)
       =(-1)^{f-1}J_{\Fq_Q}(\rho,\rho)^f;
\]
its hypotheses hold because \(\rho\), \(\rho\), and \(\rho^2\) are all
nontrivial \cite[Theorem~1.1]{HoshiKanai22}.  By
\Cref{eq:eisenstein-jacobi}, the right side is
\(-\pi_+(\mathfrak q)^f\).  Since
\(\Norm_{F/E_0}\mathfrak p=\mathfrak q^f\),
\Cref{eq:cubic-jacobi} proves
\(\gamma_2(\mathfrak p)^3=-\Lambda(\mathfrak p)\).
The local norm at infinity is a product of complex coordinates or their
conjugates.  Consequently its angular exponent at each complex place of
\(F\) is \(1\) or \(-1\), with no radial factor.  Pullback by norm is
unramified wherever \(\Lambda_{E_0}\) is unramified, so its finite conductor
is supported in \(S\).

For the quadratic sum we use one global additive character to relate the
local signs.  For \(a\in F_v^\times\), define the normalized Weil index
\[
 \mathfrak w_v(a)=
 \frac{\operatorname{WeilIndex}_{\psi_v}(a x^2)}
      {\operatorname{WeilIndex}_{\psi_v}(x^2)}.
\]
It depends only on the square class of \(a\).  Weil's product formula and
the Hilbert-symbol identity are
\begin{equation}
 \prod_v\mathfrak w_v(a)=1\quad(a\in F^\times),\qquad
 \mathfrak w_v(a)\mathfrak w_v(b)
       =(a,b)_{2,v}\mathfrak w_v(ab);
\label{eq:weil-index}
\end{equation}
see \cite[Chapter~II, \S28 and Proposition~5 in \S30]{Weil64}.
At an odd conductor-zero place, the local Gauss integral for the Weil index
\cite[Chapter~II, \S27]{Weil64} gives \(\mathfrak w_v(u)=1\) for a unit \(u\)
and
\[
 \mathfrak w_{\mathfrak p}(p^{-1})
  =q_{\mathfrak p}^{-1/2}
       \sum_{x\bmod\mathfrak p}\psi_{\mathfrak p}(x^2/p)
  =q_{\mathfrak p}^{-1/2}
       \sum_{x\bmod\mathfrak p}
          \chi_{\mathfrak p}(x)^3\psi_{\mathfrak p}(x/p)
  =\gamma_3(\mathfrak p).
\]
For the second equality, the number of square roots of \(x\) is
\(1+\chi_{\mathfrak p}(x)^3\), and the sum of a nontrivial additive
character is zero.  The last equality uses
\(\chi_{\mathfrak p}(-1)=1\).

Put \(\mathfrak w_S=\prod_{v\in S}\mathfrak w_v\).  Applying the first
identity of \Cref{eq:weil-index} to \(p^{-1}\), all exterior places other
than \(\mathfrak p\) contribute \(1\), and hence
\[
 \gamma_3(\mathfrak p)=\mathfrak w_S(p^{-1})^{-1}.
\]
On \(B\) the multiplicative defect in the second identity is
\[
 \prod_{v\in S}(a,b)_{2,v}=H_S(a,b)^3=1.
\]
Thus \([a]\mapsto\mathfrak w_S(a^{-1})^{-1}\) is a character of \(B\)
factoring through square classes.  It takes values in \(\{\pm1\}\);
equivalently Gauss inversion gives
\(\gamma_3(\mathfrak p)^2=\chi_{\mathfrak p}(-1)^3=1\).
Extend this B-character by \Cref{lem:kummer-characters}.  Since \(4\) is an
S-unit, \(\chi_\bullet(4)\) is also a finite-order Hecke character unramified
outside \(S\).  Their product, with \(\chi_\bullet(4)\) conjugated, defines
\(G\) and gives its displayed prime value.

It remains to check the sextic duplication factor.  In the finite-field
notation above put \(\phi=\chi^3\).  Completing the square and counting
square roots gives
\[
 J(\chi,\chi)
  =\overline{\chi(4)}\sum_{t\in k}\chi(1-t^2)
  =\overline{\chi(4)}J(\phi,\chi).
\]
In the second equality the part without \(\phi\) vanishes, since \(\chi\)
is nontrivial.  Applying the Gauss--Jacobi identity to both sides yields
\[
 G_1G_4=\overline{\chi(4)}G_2G_3.
\]
Together with \(G_2G_4=Q\), this gives
\[
 \gamma_1(\mathfrak p)\gamma_2(\mathfrak p)
   =\overline{\chi_{\mathfrak p}(4)}\,
      \gamma_3(\mathfrak p)\gamma_2(\mathfrak p)^3
   =-\Lambda(\mathfrak p)G(\mathfrak p).
\]

Finally use the allowed product \(c_1c_2\) to evaluate the sum at
\(\mathfrak c_1\mathfrak c_2\).  Chinese remaindering writes its additive
character as the product of the characters modulo \(\mathfrak c_i\), with
their arguments multiplied respectively by \(c_2^{-1}\) and \(c_1^{-1}\).
Changing the two residue variables contributes
\(\chi_{\mathfrak c_1}(c_2)^j\chi_{\mathfrak c_2}(c_1)^j\).
This proves \Cref{eq:gauss-crt}, including \(j\equiv0\pmod6\) with the
specified zero convention.
\end{proof}

For later use set
\[
 \vartheta=\overline\Lambda\,\overline G\,\eta.
\]
Its finite conductor is supported in \(S\).  If \(\vartheta_v\) denote its
idele-class components, define
\[
 L_S(y)=\prod_{v\in S}\vartheta_v(y_v)^{-1}.
\]
The global character is trivial on \(F^\times\), so for every exterior
ideal with allowed generator \(a\),
\begin{equation}
 \vartheta(\mathfrak a)=L_S(a),\qquad L_S(\epsilon)=1
       \quad(\epsilon\in U_S).\label{eq:vartheta-s}
\end{equation}
This is a smooth unitary character of \(F_S^\times\).  Finite-order
characters are trivial on the connected group \(\C^\times\); hence the
infinity components contributed by \(G\) and \(\eta\) are trivial.
Those of \(L_S\) are therefore purely angular of exponent \(1\) or \(-1\)
at each complex place.
